\documentclass[11pt,reqno]{amsart}

% --- packages ---
\usepackage[margin=1in]{geometry}
\usepackage{amsmath,amssymb,amsthm,mathtools}
\usepackage{graphicx}
\usepackage{booktabs}
\usepackage{placeins}
\usepackage{float}
\usepackage{enumitem}
\usepackage{microtype}
\usepackage{xcolor}
\usepackage{hyperref}
\usepackage[nameinlink,capitalize]{cleveref}
\usepackage{listings}
\usepackage{fancyvrb}
\DefineVerbatimEnvironment{leanshade}{Verbatim}{%
  fontsize=\footnotesize,
  frame=single,
  framesep=4pt,
  baselinestretch=1.05,
  xleftmargin=0.5em
}
\usepackage{fontspec}
\usepackage{csquotes}

% Body: Latin Modern (default XeTeX). Mono: DejaVu for Lean Unicode.
\setmonofont{DejaVu Sans Mono}[Scale=0.90]

\hypersetup{
  colorlinks=true,
  linkcolor=blue!50!black,
  citecolor=blue!50!black,
  urlcolor=blue!60!black,
  pdftitle={A Tether That Holds Nothing},
  pdfauthor={Benjamin Stanley Frohman},
}

% --- theorem environments ---
\theoremstyle{plain}
\newtheorem{theorem}{Theorem}[section]
\newtheorem{proposition}[theorem]{Proposition}
\newtheorem{lemma}[theorem]{Lemma}
\newtheorem{corollary}[theorem]{Corollary}

\theoremstyle{definition}
\newtheorem{definition}[theorem]{Definition}
\newtheorem{hypothesis}[theorem]{Hypothesis}

\theoremstyle{remark}
\newtheorem{remark}[theorem]{Remark}
\newtheorem{example}[theorem]{Example}

% --- macros ---
\newcommand{\T}{\mathfrak{T}}
\newcommand{\Tvel}{\mathfrak{T}_{\mathrm{vel}}}
\newcommand{\TT}{\mathbb{T}}
\newcommand{\R}{\mathbb{R}}
\newcommand{\Z}{\mathbb{Z}}
\newcommand{\enst}{\|\omega\|_{L^2}^2}
\newcommand{\Lean}[1]{\texttt{\detokenize{#1}}}
\newcommand{\dd}{\mathrm{d}}
\DeclareMathOperator{\curl}{curl}
\DeclareMathOperator{\BS}{BS}
\DeclareMathOperator{\diag}{diag}

% --- listings ---
\lstdefinestyle{lean}{
  basicstyle=\ttfamily\footnotesize,
  breaklines=true,
  breakatwhitespace=false,
  frame=single,
  framerule=0.4pt,
  rulecolor=\color{black!40},
  backgroundcolor=\color{gray!6},
  columns=fullflexible,
  keepspaces=true,
  showstringspaces=false,
  aboveskip=0.8em,
  belowskip=0.6em,
  xleftmargin=0.6em,
  xrightmargin=0.4em,
  literate=
    {ω}{{ω}}1
    {δ}{{δ}}1
    {κ}{{κ}}1
    {Π}{{Π}}1
    {ℝ}{{ℝ}}1
    {∫}{{∫}}1
    {‖}{{‖}}1
    {²}{{²}}1
    {∂}{{∂}}1
    {λ}{{λ}}1
    {·}{{·}}1
    {≤}{{≤}}1
    {→}{{→}}1
    {∀}{{∀}}1
    {∃}{{∃}}1
    {≠}{{≠}}1
    {⟨}{{⟨}}1
    {⟩}{{⟩}}1
    {⟪}{{⟨}}1
    {⟫}{{⟩}}1,
}
\lstset{style=lean}

% --- title ---
\title[A Tether That Holds Nothing]{%
  A Tether That Holds Nothing:\\
  Energy-Null Metriplectic Corrections\\
  and the Vanishing Enstrophy Contribution on $\TT^3$}

\author{Benjamin Stanley Frohman}
\address{Independent researcher}
\email{frohmanbenjamin@gmail.com}
\thanks{ORCID: 0009-0006-7068-3718}
\date{4 November 2026}

\subjclass[2020]{35Q30, 76D05, 37K65}
\keywords{Navier--Stokes equations, metriplectic systems, coadjoint orbits,
enstrophy, Beale--Kato--Majda, geometric tether, formal Lean verification}

\begin{document}
\emergencystretch=2em
\begin{abstract}
We record the Frohmanian geometric tether---a symmetric, negative-semidefinite
quadratic correction to the Lie--Poisson bracket on the vorticity coadjoint
orbit of ideal incompressible fluid---in both its stream-function ($\psi$) and
velocity-slot forms. The velocity-slot kernel $\Tvel$ annihilates the kinetic
energy by construction (condition~C2), restoring consistency of the metriplectic
sum with the classical Euler equations. That consistency, however, is precisely
the obstruction: by symmetry the tether contributes nothing to the enstrophy
balance, so no cubic damping of enstrophy appears. We state the open analytic
hypothesis \Lean{FrohmanianTetherRate}, distinguish its inequivalent differential
forms, and delimit the scope of the associated product theorem: unconditional
only for a discontinuous step field, zero data, and constant flow. No claim is
made on the Clay Millennium problem for the three-dimensional Navier--Stokes
equations.
\end{abstract}

\maketitle

% =====================================================================
\section{Introduction}\label{sec:intro}
% =====================================================================

The geometric structure of ideal incompressible fluid on a compact Riemannian
manifold is classical: the Euler equations are Hamiltonian on the coadjoint
orbits of the volume-preserving diffeomorphism group, with the Lie--Poisson
bracket of Arnold~\cite{Arnold1966,ArnoldKhesin} and kinetic energy as
Hamiltonian. Dissipative corrections that preserve energy while producing
entropy growth---metriplectic systems in the sense of
Morrison~\cite{Morrison1986} and Grmela--\"{O}ttinger~\cite{GrmelaOttinger1997}---offer
a natural language in which to ask whether a structurally minimal correction
could supply the missing a-priori bound on vorticity.

The \emph{Frohmanian geometric tether} is one such candidate: a symmetric
bilinear form $\T(F,G)$ on functionals of vorticity, weighted by global
enstrophy and projected orthogonal to the Biot--Savart velocity. Its design
goals are three conditions conventionally labelled C1--C3:
\begin{enumerate}[label=(C\arabic*),leftmargin=2.2em]
  \item negative-semidefiniteness of the quadratic form $\T(F,F)$;
  \item annihilation of the kinetic energy, $\T(H,\cdot)=0$;
  \item compatibility with the coadjoint (vorticity) structure.
\end{enumerate}
Condition~(C2) ensures that the tethered dynamics leave the ideal Euler
equations unchanged whenever the dissipative (viscous) term is switched off.
The present note records the precise definitions used in the accompanying Lean
development \Lean{frohmanian-tether-sdiff}, the quadratic geometry of the
kernel, a C2 failure of the $\psi$-slot form, the velocity-slot repair, and the
analytic obstruction that follows once C2 holds.

The conclusion is negative in a precise sense. The very property that makes
the tether consistent with Euler---energy nullity---forces, by symmetry, a
vanishing contribution to the enstrophy balance. No damping of the form
$-\kappa'M^2$ is produced. Closing a differential inequality for enstrophy
along physical Navier--Stokes solutions therefore remains an open analytic
step, named \Lean{FrohmanianTetherRate} (also \Lean{hDI}). The Clay question
of global regularity for smooth finite-energy data on $\TT^3$
\cite{Fefferman2000} is untouched.

% =====================================================================
\section{Notation and classical identities}\label{sec:notation}
% =====================================================================

Let $\TT^3=\R^3/(2\pi\Z)^3$ be the flat three-torus. Write $u$ for a
divergence-free velocity field, $\omega=\curl u$ for its vorticity, and
$u=\BS(\omega)$ for the Biot--Savart reconstruction (the unique
divergence-free, mean-zero solution of $\curl u=\omega$). Functionals
$F,G$ of vorticity take values in $\R$; their functional derivatives
$\delta F/\delta\omega$ are identified with stream-function-like fields.
The kinetic energy and (unsigned) enstrophy are
\begin{equation}\label{eq:HM}
  H(u)=\tfrac12\int_{\TT^3}|u|^2\,\dd x,\qquad
  M(\omega)=\tfrac12\int_{\TT^3}|\omega|^2\,\dd x.
\end{equation}
(The factor $\tfrac12$ is conventional; rescaling $M$ rescales constants in
any differential inequality but does not alter its shape.) The classical
enstrophy identity for a smooth solution of the incompressible
Navier--Stokes equations with viscosity $\nu>0$ reads
\begin{equation}\label{eq:enstrophy}
  \frac{\dd}{\dd t}M+\nu\|\nabla\omega\|_{L^2}^2
  =\int_{\TT^3}\bigl\langle\omega,(\omega\cdot\nabla)u\bigr\rangle\,\dd x.
\end{equation}
The right-hand side is the vortex-stretching term. In three dimensions
enstrophy is \emph{not} a Casimir of the Euler--Poisson bracket: the
classical contribution $\{M,H\}_{\mathrm{LP}}$ is precisely the stretching
integral and is nonzero in general. (Helicity is a true Casimir; we do not
use it here.)

The Beale--Kato--Majda criterion~\cite{BKM1984} states that a smooth solution
on $[0,T)$ extends past $T$ whenever
\begin{equation}\label{eq:BKM}
  \int_0^T\|\omega(t)\|_{L^\infty}\,\dd t<\infty.
\end{equation}
Thus a time-integrated bound on a critical norm is the realistic analytic
target; a pointwise cubic differential inequality is stronger than what BKM
requires.

% =====================================================================
\section{The geometric tether}\label{sec:tether}
% =====================================================================

\subsection{The \texorpdfstring{$\psi$}{psi}-slot kernel}

\begin{definition}[TetherKernel]\label{def:TetherKernel}
Fix $\kappa>0$. For vorticity $\omega$ on a coadjoint orbit and functionals
$F,G$, the \emph{$\psi$-slot tether kernel} is
\begin{equation}\label{eq:T}
  \T(F,G)(\omega)
  =-\kappa\,\enst
    \int_{\TT^3}
      \Bigl\langle
        \Pi_u\bigl(\tfrac{\delta F}{\delta\omega}\bigr),\,
        \Pi_u\bigl(\tfrac{\delta G}{\delta\omega}\bigr)
      \Bigr\rangle\,\dd x,
\end{equation}
where $u=\BS(\omega)$ and $\Pi_u$ denotes the pointwise orthogonal projection
onto the plane perpendicular to $u(x)$ (with the convention $\Pi_0=\mathrm{Id}$).
In the Lean development the constant is fixed at $\kappa=\tfrac32$. The global
scalar weight $\enst$ replaced an earlier pointwise density $|\omega(x)|^2$
(author decision of 3~November~2026); the two objects are inequivalent.
\end{definition}

The Lean realisation in \Lean{frohmanian-tether-sdiff} is:

\begin{leanshade}
/-! ## The Quadratic Metric Correction (the Tether) -/
/-- Canonical Frohmanian tether kernel
`B(F,G)(ω) = -κ ‖ω‖² ∫ (Π_u (δF/δω) · Π_u (δG/δω)) dλ`, with the GLOBAL
enstrophy `‖ω‖² = Enstrophy ω = ∫ |ω|²` outside the integral. -/
@[expose] public noncomputable def TetherKernel
    (ω : CoadjointOrbit) (F G : Functional) : ℝ :=
  -kappa * Enstrophy ω * ∫ x,
      inner ℝ
        (Pi_u (velocity_from_vorticity ω) (FunctionalDerivative F ω) x)
        (Pi_u (velocity_from_vorticity ω) (FunctionalDerivative G ω) x)
    ∂volume
\end{leanshade}

\begin{leanshade}
public noncomputable def TetheredBracket
    (F G : Functional) (ω : CoadjointOrbit) : ℝ :=
  ClassicalBracket F G ω + TetherKernel ω F G
\end{leanshade}

\begin{proposition}[Quadratic form]\label{prop:quadratic}
The form $\T$ is symmetric in $(F,G)$. Its diagonal satisfies
\begin{equation}\label{eq:quad}
  \T(F,F)=-\kappa\,\enst\int_{\TT^3}\bigl|\Pi_u\tfrac{\delta F}{\delta\omega}\bigr|^2\,\dd x\le 0.
\end{equation}
In particular $\T(F,F)=0$ whenever $\delta F/\delta\omega$ is pointwise parallel
to $u$. Lean identifier:
{\footnotesize\Lean{tetherKernel_quadratic_form}}.
\end{proposition}

The tethered bracket
$\{F,G\}_{\T}=\{F,G\}_{\mathrm{LP}}+\T(F,G)$ is therefore metriplectic rather
than Poisson: the Lie--Poisson summand is antisymmetric (and Jacobi, cited but
not formalized in the present development), while the correction is symmetric.
The identity $\{F,F\}_{\T}=\T(F,F)\le 0$ encodes dissipation along level sets of
$F$. Canonicity of the kernel is proved only inside a restricted class of
quadratic corrections and under the explicit hypotheses of
{\footnotesize\Lean{uniqueness_of_minimal_tether}}; it is not an unconditional
uniqueness statement.

\FloatBarrier

\subsection{Schematic geometry of the quadratic form}

All figures in this note are schematic illustrations of the sign structure of
$\T(F,F)$. They are \emph{not} computed fluid fields.

\begin{figure}[ht]
\centering
\includegraphics[width=0.48\textwidth]{tether_3d_clean.png}\hfill
\includegraphics[width=0.48\textwidth]{tether_quadratic.png}
\caption{Schematic. Left and right: the diagonal $\T(F,F)$ over a
two-dimensional slice of argument space. The surface is identically zero along
the component parallel to $u$ (the kinetic-energy direction) and
negative-quadratic in the orthogonal component. Captions refer to
$\T(F,F)$, not to the off-diagonal $\T(F,G)$.}
\label{fig:slice}
\end{figure}

\begin{figure}[ht]
\centering
\includegraphics[width=0.48\textwidth]{tether_schematic.png}\hfill
\includegraphics[width=0.48\textwidth]{tether_1d.png}
\caption{Schematic. Left: trough geometry with energy and projected axes.
Right: one-dimensional profile $\T(F,F)\propto-|\Pi_u\delta F|^2$.}
\label{fig:trough}
\end{figure}

\begin{figure}[ht]
\centering
\includegraphics[width=0.56\textwidth]{tether_3d_corrected.png}\\[0.5em]
\includegraphics[width=0.46\textwidth]{tether_schematic2.png}
\caption{Schematic. Top: bowl $-(y^2+z^2)$ over the projected plane
$(\Pi_u)_1\times(\Pi_u)_2$. The kinetic-energy direction is orthogonal to this
plane and sits at the unique point where $\T=0$ (the origin). Earlier drafts
incorrectly drew the energy line along a projected axis; see
\cref{app:drafts}. Bottom: a coadjoint orbit (closed curve) with sample
projected vectors $\Pi_u\delta F$ (arrows). Illustration only.}
\label{fig:bowl}\label{fig:orbit}
\end{figure}

\FloatBarrier

% =====================================================================
\section{The velocity-slot repair and condition C2}\label{sec:C2}
% =====================================================================

\subsection{Failure of C2 in the \texorpdfstring{$\psi$}{psi}-slot}

Condition~(C2) asks that $\T(H,\cdot)=0$ for the kinetic energy~$H$. In the
$\psi$-slot form~\eqref{eq:T} one projects $\delta H/\delta\omega=\psi$, the
stream function. In general $\psi\nparallel u$, so
$\Pi_u\psi\not\equiv 0$. On 30~October~2026 this was confirmed in Lean on a
two-mode shear flow (commit~\texttt{a52c6af}): the $\psi$-slot kernel fails C2
on that example. The global enstrophy weight, being a scalar factor, cannot
repair a pointwise nonvanishing integrand.

\subsection{The velocity-slot kernel}

The repair is to project on the velocity side of the functional derivative.

\begin{definition}[TetherKernelVel]\label{def:TetherKernelVel}
The \emph{velocity-slot tether kernel} is
\begin{equation}\label{eq:Tvel}
  \Tvel(F,G)(\omega)
  =-\kappa\,\enst
    \int_{\TT^3}
      \Bigl\langle
        \Pi_u\Bigl(\curl\tfrac{\delta F}{\delta\omega}\Bigr),\,
        \Pi_u\Bigl(\curl\tfrac{\delta G}{\delta\omega}\Bigr)
      \Bigr\rangle\,\dd x.
\end{equation}
\end{definition}

\begin{leanshade}
public noncomputable def TetherKernelVel
    (ω : CoadjointOrbit) (F G : Functional) : ℝ :=
  -kappa * Enstrophy ω * ∫ x,
      inner ℝ
        (Pi_u (velocity_from_vorticity ω)
              (curl (FunctionalDerivative F ω)) x)
        (Pi_u (velocity_from_vorticity ω)
              (curl (FunctionalDerivative G ω)) x)
    ∂volume
\end{leanshade}

\begin{proposition}[C2 for the velocity slot]\label{prop:C2}
Assume $\curl(\delta H/\delta\omega)=u=\BS(\omega)$. Then
$\Pi_u u=0$ pointwise, hence $\Tvel(H,G)=0$ for every functional~$G$.
\end{proposition}

\begin{proof}
By the standing identification of the functional derivative of kinetic energy,
$\delta H/\delta\omega=\psi$ with $\curl\psi=u$. The integrand of~\eqref{eq:Tvel}
at $F=H$ therefore contains the factor $\Pi_u u$, which vanishes at every
$x\in\TT^3$ by definition of $\Pi_u$. The scalar prefactor $-\kappa\enst$ is
immaterial. Thus $\Tvel(H,G)=0$.
\end{proof}

\begin{remark}[Lean status]\label{rem:lean-status}
The argument of \cref{prop:C2} is elementary once the identification
$\curl(\delta H/\delta\omega)=u$ is available in the development. At the time of
writing \Lean{tetherKernelVel_C2} is kernel-checked in Lean~4 (0 \texttt{sorry}; axioms
\texttt{propext}, \texttt{Classical.choice}, \texttt{Quot.sound} only; local commit \texttt{7e3cf78}),
\emph{conditional} on the named open hypothesis \Lean{KineticEnergyCurlIdentification},
\[\curl(\delta H/\delta\omega)=\mathrm{velocity\_from\_vorticity}\,\omega,\]
together with $\operatorname{div}u=0$ and finite energy. Because $\Pi_u$ in the
development is the $L^2$ Gram--Schmidt projector rather than a pointwise one,
$\Pi_u u\,u=0$ is proved almost everywhere under those standing hypotheses
(\Lean{Pi_u_self_ae_zero}). Discharging the curl identification requires a G\^ateaux
derivative and Biot--Savart calculus not yet present in the development.
\end{remark}

% =====================================================================
\section{The C2 obstruction to enstrophy damping}\label{sec:obstruction}
% =====================================================================

Suppose C2 holds in the form $\T(H,G)=0$ for every second slot~$G$ (as it does
for $\Tvel$ under the hypotheses of \cref{prop:C2}). Taking $G=M$, the
enstrophy functional, yields $\T(H,M)=0$. Symmetry of $\T$ then gives
\begin{equation}\label{eq:symmetry}
  \T(M,H)=\T(H,M)=0.
\end{equation}
Along the tethered flow the evolution of enstrophy is therefore
\begin{equation}\label{eq:dM}
  \frac{\dd M}{\dd t}
  =\{M,H\}_{\mathrm{LP}}+\T(M,H)
  =\{M,H\}_{\mathrm{LP}}.
\end{equation}
The tether contributes nothing. In particular no term of the shape
$-\kappa'M^2$ (or $-\kappa''M^3$) is generated by the geometric correction. The
classical stretching contribution $\{M,H\}_{\mathrm{LP}}$ remains exactly as in
the untethered Euler or Navier--Stokes dynamics.

\begin{remark}[Enstrophy is not a 3D Casimir]\label{rem:not-casimir}
The vanishing~\eqref{eq:symmetry} should not be confused with the Casimir
property. In three dimensions, enstrophy is not a Casimir of the Euler
bracket: $\{M,H\}_{\mathrm{LP}}$ equals the vortex-stretching integral and is
nonzero for generic data. The tether simply fails to modify that classical
balance. Replacing ``enstrophy'' by a true Casimir such as helicity yields the
same formal conclusion---the tether adds nothing---but helicity does not
control regularity.
\end{remark}

\begin{remark}[Projection gap]\label{rem:projection-gap}
Even an asymmetric tether designed to produce damping would confront a
projection gap. The only estimate available from the definition is the upper
bound
\[
  \int_{\TT^3}|\Pi_u\omega|^2\,\dd x\le\|\omega\|_{L^2}^2.
\]
A coercive lower bound on $|\Pi_u\omega|$ is what damping requires, and it
fails wherever $\omega\parallel u$---precisely the Beltrami-type regions in which
vortex stretching is weakest. No such lower bound follows from the present
geometric construction.
\end{remark}

Thus two routes have been examined.
\begin{enumerate}[leftmargin=1.5em]
  \item \emph{Inside the bracket} (metriplectic tether). The C2/C3 obstruction
    prevents the symmetric tether from producing cubic damping. No modification
    of the present form removes the obstruction while preserving both energy
    conservation and the coadjoint structure.
  \item \emph{Outside the bracket} (a-posteriori estimate). Starting from a
    smooth solution of the \emph{unmodified} 3D Navier--Stokes equations, one
    may form a Lyapunov quantity from its vorticity and attempt a differential
    inequality by classical identities. This avoids the C2 conflict, but the
    resulting inequality does not close for arbitrary data.
\end{enumerate}
Until an open rate hypothesis is discharged by a genuine estimate, neither
route yields global regularity for every smooth finite-energy initial datum.

% =====================================================================
\section{The open hypothesis \texorpdfstring{\Lean{FrohmanianTetherRate}}{FrohmanianTetherRate}}\label{sec:rate}
% =====================================================================

\begin{hypothesis}[FrohmanianTetherRate / hDI]\label{hyp:rate}
There exist constants $C>0$ and $\kappa'>0$ such that every smooth solution of
the 3D incompressible Navier--Stokes equations on $\TT^3$ satisfies a
differential inequality of one of the following shapes for its enstrophy
(or a supremum-norm proxy) $M(t)$:
\begin{align}
  M'(t)&\le C\,M(t)^2-\kappa''\,M(t)^3,\label{eq:rate-cubic}\\
  M'(t)&\le C\,M(t)^{3/2}-\kappa'\,M(t)^2.\label{eq:rate-threehalf}
\end{align}
\end{hypothesis}

\begin{remark}[Inequivalence]\label{rem:inequiv}
The two forms~\eqref{eq:rate-cubic} and~\eqref{eq:rate-threehalf} are
\emph{not} equivalent: they assert different powers and neither implies the
other. A single official form must be fixed; referring to both as ``equivalent''
is a paper-to-code mismatch. The factor $\tfrac12$ in~\eqref{eq:HM} merely
rescales constants and cannot convert an unproved inequality into a proved one.
\end{remark}

The hypothesis is taken as an explicit argument by the conditional regularity
theorems {\footnotesize\Lean{frohmanian_tether_theorem}} and
{\footnotesize\Lean{global_regularity_for_NS}}.
It is \emph{not} derived from the classical identity~\eqref{eq:enstrophy}, nor
from the geometric tether. Earlier attempts to obtain a pointwise cubic bound
of this shape were withdrawn on grounds of scale inconsistency and the C2
obstruction of~\cref{sec:obstruction}.

Known obstructions to a direct derivation include:
\begin{itemize}[leftmargin=1.2em]
  \item \emph{Scale inconsistency.} A sink $-\kappa M^4$, or $-\kappa'M^2$ with
    dimensionless $\kappa'$, does not match the parabolic scaling of
    Navier--Stokes.
  \item \emph{The C2 obstruction.} The symmetric tether contributes zero to the
    enstrophy balance~\eqref{eq:dM}.
  \item \emph{The projection gap.} Only an upper bound on $\Pi_u\omega$ is
    available; cf.~\cref{rem:projection-gap}.
\end{itemize}
By Beale--Kato--Majda~\eqref{eq:BKM}, a time-integrated bound on
$\|\omega\|_\infty$ already yields continuation. The differential
inequalities~\eqref{eq:rate-cubic}--\eqref{eq:rate-threehalf} are therefore
stronger than necessary; a weaker integrated form would suffice for the
regularity application and may be a more realistic target.

Concrete next steps that do not claim a solution:
\begin{enumerate}[leftmargin=1.5em]
  \item Extract the exact classical stretching integral and the precise
    H\"{o}lder/Sobolev constants already present in the Lean files.
  \item Test candidate scale-invariant replacements (for instance
    $-\kappa\nu M^2\cdot(\text{dimensionless quantity})$) against known exact
    solutions and scaling counter-examples.
  \item Record any partial estimate that holds under an extra hypothesis---e.g.\
    a bound on the direction of vorticity in the sense of
    Constantin--Fefferman~\cite{CF1993}---as a named conditional lemma.
\end{enumerate}

% =====================================================================
\section{Scope of the product theorem}\label{sec:scope}
% =====================================================================

The Lean declaration
\begin{quote}
\raggedright\footnotesize\ttfamily
frohmanian\_metriplectic\_system\_canonical\_uniqueness\_and\_global\_\-regularity\_product\_theorem
\end{quote}
is assumption-free as a formal object, but its regularity conjunct covers only
three regimes:
\begin{enumerate}[label=(\roman*),leftmargin=1.8em]
  \item a discontinuous step field that is not a Navier--Stokes solution;
  \item identically zero data;
  \item constant (trivial) flow.
\end{enumerate}
Every theorem in the development that addresses general smooth finite-energy
data takes one of the named open hypotheses
{\footnotesize\Lean{hRiccati}},
{\footnotesize\Lean{FrohmanianTetherBound}}, or
{\footnotesize\Lean{TetherCoefficientBounded}}
as an explicit argument. Consequently:

\begin{theorem}[Scope; informal]\label{thm:scope}
There is no unconditional global-regularity theorem for arbitrary smooth
finite-energy initial data of 3D Navier--Stokes on $\TT^3$ in the present
development. The Clay Millennium problem~\cite{Fefferman2000} is not solved
here.
\end{theorem}

The honest forward path is therefore to drop the tether from the dynamics and
treat any candidate bound as an a-posteriori estimate on true Navier--Stokes
solutions, labelled explicitly as conditional on the rate or absorbing-bound
hypotheses above.

% =====================================================================
\section*{Acknowledgements}
% =====================================================================

The accompanying formal development is the Lean repository
\Lean{frohmanian-tether-sdiff}. Schematic figures were generated from the
quadratic form~\eqref{eq:quad} for illustration only.

% =====================================================================
\appendix
\section{Earlier drafts of the three-dimensional figure}\label{app:drafts}
% =====================================================================

The two plots below are the author's original script outputs, retained for the
record. In both, a black ``energy direction'' line was drawn along a projected
axis, incorrectly placing the kinetic-energy direction inside the projected
plane. The corrected \cref{fig:bowl} shows the energy direction orthogonal to
the projected plane, at the unique point of the figure where $\T=0$.

\begin{figure}[ht]
\centering
\includegraphics[width=0.48\textwidth]{tether_3d.png}\hfill
\includegraphics[width=0.48\textwidth]{tether_3d_v2.png}
\caption{Earlier drafts (author's original script outputs). Schematic only.
The energy-direction line coincides with a projected axis---an error corrected
in \cref{fig:bowl}.}
\label{fig:drafts}
\end{figure}

% =====================================================================
\begin{thebibliography}{99}

\bibitem[Arn66]{Arnold1966}
V.~I. Arnold,
\newblock Sur la g\'{e}om\'{e}trie diff\'{e}rentielle des groupes de Lie de
  dimension infinie et ses applications \`{a} l'hydrodynamique des fluides
  parfaits,
\newblock \emph{Ann.\ Inst.\ Fourier (Grenoble)} \textbf{16} (1966), 319--361.

\bibitem[AK98]{ArnoldKhesin}
V.~I. Arnold and B.~A. Khesin,
\newblock \emph{Topological Methods in Hydrodynamics},
\newblock Appl.\ Math.\ Sci.\ 125, Springer, New York, 1998.

\bibitem[BKM84]{BKM1984}
J.~T. Beale, T. Kato, and A. Majda,
\newblock Remarks on the breakdown of smooth solutions for the 3-D Euler
  equations,
\newblock \emph{Comm.\ Math.\ Phys.} \textbf{94} (1984), 61--66.

\bibitem[CF93]{CF1993}
P. Constantin and C. Fefferman,
\newblock Direction of vorticity and the problem of global regularity for the
  Navier--Stokes equations,
\newblock \emph{Indiana Univ.\ Math.\ J.} \textbf{42} (1993), 775--789.

\bibitem[Fef00]{Fefferman2000}
C.~L. Fefferman,
\newblock Existence and smoothness of the Navier--Stokes equation,
\newblock in \emph{The Millennium Prize Problems}, Clay Mathematics Institute,
  Cambridge, MA, 2000.

\bibitem[GO97]{GrmelaOttinger1997}
M. Grmela and H.~C. \"{O}ttinger,
\newblock Dynamics and thermodynamics of complex fluids. I.~Development of a
  general formalism,
\newblock \emph{Phys.\ Rev.\ E} \textbf{56} (1997), 6620--6632.

\bibitem[Mor86]{Morrison1986}
P.~J. Morrison,
\newblock A paradigm for joined Hamiltonian and dissipative systems,
\newblock \emph{Physica D} \textbf{18} (1986), 410--419.

\end{thebibliography}
\end{document}
